\section{Continuum substrate model}
\label{sec:continuum-model}

\subsection{Ontology and kinematics (time as parameter)}

Let $\Omega\subset\R^3$ be the intrinsic (material) brane domain.
We use either a periodic cell $\Omega=\T^3$ or a bounded cube with clamped boundary points.
Time $t\in\R$ is an external evolution parameter, not an additional coordinate of the brane.

The fundamental field is the embedding map
\begin{equation}
  \vecX(x,t)\in\R^4,
  \qquad x=(x^1,x^2,x^3)\in\Omega,\quad t\in\R,
  \label{eq:fundamental-field}
\end{equation}
whose components $X^A(x,t)$ ($A=1,\dots,4$) give the brane configuration in the ambient space.

\subsection{Induced metric and isotropic reference metric}

The embedding induces a (time-dependent) Riemannian metric on the brane:
\begin{equation}
  g_{ij}(x,t) := \partial_i \vecX(x,t)\cdot \partial_j \vecX(x,t),
  \qquad i,j=1,2,3.
  \label{eq:induced-metric}
\end{equation}
This definition is invariant under rigid motions $\vecX\mapsto Q\vecX+a$ with $Q\in O(4)$ and $a\in\R^4$.
Rotational invariance within the brane is obtained by constructing the stored energy density from invariants of $g$
relative to an isotropic reference metric.

We choose a homogeneous, isotropic reference metric
\begin{equation}
  g^0_{ij}=\ell_0^2\,\delta_{ij},
  \label{eq:reference-metric}
\end{equation}
where $\ell_0>0$ is the continuum analogue of a discrete spring rest length.
Define the mixed relative metric
\begin{equation}
  \widehat g := (g^0)^{-1}g.
  \label{eq:relative-metric}
\end{equation}
Any isotropic hyperelastic energy density can be written as $W=W(\widehat g)$, equivalently as a function of the
principal stretches or invariants of $\widehat g$.

\paragraph{Continuum analogue of pre-tension (the $\alpha$ knob).}
The reference metric $g^0$ encodes the stress-free geometric scale.
If boundary conditions enforce a held ground-state scale $h_\star$ (uniform), then
\begin{equation}
  \alpha:=\frac{\ell_0}{h_\star}\in(0,1]
  \label{eq:alpha-def}
\end{equation}
measures the mismatch between stress-free scale and held scale.
For $\alpha<1$ the brane is uniformly pre-stressed (continuum ``pre-tension''). This mismatch controls how strongly
transverse gradients backreact on in-brane stresses through the geometric nonlinearity of $g_{ij}(\vecX)$.

\subsection{Constitutive law: St.\ Venant--Kirchhoff (Green--Lagrange strain)}

As a concrete baseline we use the St.\ Venant--Kirchhoff (StVK) energy expressed in terms of the
Green--Lagrange strain tensor
\begin{equation}
  E := \tfrac12\big(\widehat g - I\big)
  \qquad\text{(with $I$ the $3\times 3$ identity)}.
  \label{eq:green-lagrange-strain}
\end{equation}
The StVK stored energy density is
\begin{equation}
  W_{\text{StVK}}(E)
  = \mu\,\mathrm{tr}(E^2)+\frac{\lambda}{2}\,\big(\mathrm{tr}E\big)^2,
  \label{eq:stvk}
\end{equation}
with Lam\'e parameters $\mu>0$ and $\lambda$ (units: Pa). In terms of Poisson ratio $\nu\in(0,1/2)$,
\begin{equation}
\nu=\frac{\lambda}{2(\lambda+\mu)},
\qquad
\lambda=\frac{2\mu\nu}{1-2\nu}.
\label{eq:poisson-lame}
\end{equation}
This choice is not claimed unique; it is used because it is isotropic, analytic, and captures the
geometric nonlinearity through $g_{ij}(\vecX)$.

\subsection{Lagrangian, action, and equations of motion}

Let $\rho_m>0$ be the brane mass density (kg/m$^3$). The Lagrangian density is
\begin{equation}
  \Lag(\vecX,\partial_t\vecX,\nabla\vecX)
  =\frac{\rho_m}{2}\,|\partial_t\vecX|^2
  -W\!\left(g(\vecX);g^0\right).
  \label{eq:lagrangian-density}
\end{equation}
The action is
\begin{equation}
  S[\vecX] = \int_{\R}\dd t \int_{\Omega} \dd^3 x\;\Lag.
  \label{eq:action}
\end{equation}

Define the first Piola stress (as a momentum conjugate to $\partial_i X^A$)
\begin{equation}
  P_A^{\;i} := \frac{\partial W}{\partial(\partial_i X^A)}.
  \label{eq:piola-def}
\end{equation}
For energies depending on $g_{ij}$ one can write $P_A^{\;i}=2S^{ij}\partial_j X^A$ with a symmetric
second Piola-type tensor $S^{ij}=\partial W/\partial g_{ij}$.
The Euler--Lagrange equations take divergence form
\begin{equation}
  \rho_m\,\partial_{tt} X^A = \partial_i P_A^{\;i},
  \qquad A=1,\dots,4,
  \label{eq:eom}
\end{equation}
showing explicitly that rotational invariance within the brane is built in through the invariant construction
of $W(g;g^0)$.

\subsection{Spherical material coordinates (coordinate change only)}
\label{subsec:spherical-material-coords}

\paragraph{Important: this is only a coordinate change.}
All dynamical statements remain those of the coordinate-free formulation
$\vecX:\Omega\times\R\to\R^4$ with induced metric $g_{ij}=\partial_i\vecX\cdot\partial_j\vecX$.
We introduce spherical coordinates on the \emph{intrinsic/material} domain as a convenient chart for localized
excitations; this does \emph{not} constrain the embedding $\vecX$.

Let $(r,\theta,\varphi)$ be the standard spherical chart on (a subset of) $\R^3$, with Jacobian
$J(r,\theta,\varphi)=r^2\sin\theta$.
Writing the intrinsic coordinate as $q^a\in\{r,\theta,\varphi\}$ and denoting derivatives in this chart by $\partial_a$,
the equations of motion become
\begin{equation}
  \rho_m\,\partial_{tt}X^A
  =\frac{1}{J(q)}\,\partial_a\!\Big(J(q)\,P_A^{\;a}\Big),
  \qquad A=1,\dots,4,
  \label{eq:eom-spherical}
\end{equation}
which is exactly the Cartesian divergence form \eqref{eq:eom} written with the standard Jacobian factor.

\subsection{Linearization and polarized branches}

Linearize around a homogeneous background embedding $\vecX_\star(x)=(x,0)$ and small displacements
$\vecu(x,t)=\vecX(x,t)-\vecX_\star(x)$.
Under standard regularity assumptions on $W$, one obtains a linear wave system with (in general) multiple
polarization branches, including transverse and longitudinal modes.
For later use, we emphasize two generic features:
(i) different branches may have different characteristic speeds (and may become dispersive at shorter wavelengths),
and (ii) for $\alpha<1$ the pre-stress couples transverse excursions (in the embedding) back into in-brane stresses
through the induced metric \eqref{eq:induced-metric}.
These branches provide the substrate carriers for the geometric-phase and soliton constructions developed below.
